\section{관련 연구}
\label{sec:related}

\begin{revision}

\subsection{추론 기반 주행과 추론 충실성}
\label{sec:related-reasoning}

DriveLM은 지각ㆍ예측ㆍ계획의 질의응답을 그래프로 연결하고, DriveVLM은 장면 설명ㆍ분석ㆍ계층적 계획을 결합한다~\cite{sima2024drivelm,tian2024drivevlm}. Alpamayo-R1은 CoC 추론과 궤적 예측을 결합하고, 지도 학습과 강화 학습으로 추론과 행동의 정렬을 높인다~\cite{wang2025alpamayo}. 이 연구들은 주행 판단을 학습 가능한 언어 표현으로 만드는 데 초점을 둔다.

생성된 추론이 실제 판단을 반영하는지는 별개의 문제다. Turpin 등은 사고 과정 설명이 답에 영향을 준 요인을 빠뜨릴 수 있음을 보였고~\cite{turpin2023}, Lanham 등은 예측이 제시된 추론에 실제로 의존하는지를 측정했다~\cite{lanham2023}. 주행 분야에서 Drive the Thoughts는 VLA 추론과 예측 궤적의 일관성을 실행 중에 감시하는 모니터를 제안했다~\cite{yu2026drivethoughts}. 언어 근거와 장면ㆍ궤적의 충실성 및 시각 교란에 대한 반응을 평가하거나~\cite{mayumu2026faithfulness}, 모순ㆍ환각 지표와 CoT 개입을 함께 사용하는 연구도 있다~\cite{nguyen2026vladrivebench}.

이 연구들은 한 장면에서 설명이 행동과 맞는지, 또는 설명을 바꾸면 행동이 달라지는지를 묻는다. 본 연구는 모델을 다시 실행하거나 개입하지 않고, 저장된 연속 설명에서 앞선 의무와 뒤따르는 해제 근거를 연결한다. 따라서 생성 모델의 인과적 충실성을 검증하는 방법을 대신하지 않는다.

\subsection{교통 규칙의 형식화와 검증}
\label{sec:related-rules}

허용되는 주행 행동을 형식적으로 표현하려는 연구가 있다. Esterle 등은 교통 규칙을 선형 시간 논리(LTL)로 기술했고~\cite{esterle2020}, Maierhofer 등은 고속도로 규칙을 시간 논리로 형식화하여 준수 여부를 평가했다~\cite{maierhofer2020}. RSS는 안전 행동과 책임의 수학적 조건을 정의한다~\cite{shalev2017rss}. 형식화된 규칙 자체를 검증하는 연구도 있다. Wan 등은 등가성 검증과 모델 검사로 디지털 교통 규칙의 의미 일관성과 정확성을 확인했다~\cite{wan2024}.

이 연구들의 검사 대상은 사람이 정한 규칙이다. 본 연구의 대상은 모델이 자동 생성한 주석에 암묵적으로 담긴 의무이다. 규칙이 미리 주어지지 않으므로, 의무가 언제 시작되고 무엇으로 해제되는지를 주석에서 읽어 내야 한다.

\subsection{자연어 요구의 형식화}
\label{sec:related-nl}

자연어 요구를 형식 명세로 옮기는 연구가 있다. NL2TL은 원자 명제를 추상화한 뒤 요구 문장을 시간 논리식으로 변환하도록 학습한다~\cite{chen2023nl2tl}. nl2spec은 언어 구절과 하위 논리식의 대응을 보여 주어 사용자가 모호하거나 잘못된 번역을 고칠 수 있게 한다~\cite{cosler2023nl2spec}. 이 연구들의 주된 출력은 개별 요구 문장에 대한 형식 명세다.

연속 CoC에서는 한 문장이 만든 의무를 다른 문장이 해제하며, 해제 증거가 아예 없을 수도 있다. 본 연구는 원문 구절을 의무ㆍ해제 조건ㆍ증거에 연결해 남긴다는 점에서 nl2spec의 대응 추적과 목적이 비슷하다. 다만 범용 번역기를 제시하지 않으며, 본 평가의 프레임은 설계자가 작성했다.

\subsection{실행 중 검증과 모델 검사}
\label{sec:related-rv}

실행 중 검증(runtime verification)은 관측된 실행 경로가 주어진 명세를 만족하는지 검사한다~\cite{leucker2009runtime}. LTL/TLTL 모니터링은 유한 접두 경로에서 참ㆍ거짓ㆍ미결정을 구분하고 최초로 확정되는 시점을 다룬다~\cite{bauer2011runtime}. 따라서 이전 상태의 기억, 세 값의 판정과 최초 위반 추적 자체는 기존 기술이다. 본 논문의 미상 증거와 그에 따른 근거 부족 판정은 명시된 해제 증거가 없다는 뜻이며, 미래 연장 전체에 대한 LTL 세 값 의미론을 구현했다는 뜻은 아니다.

UPPAAL은 시간 오토마타의 상태와 도달 속성을 검사한다~\cite{behrmann2004uppaal,uppaalqueries}. 본 연구의 모델은 사건 배열을 순서대로 처리하는 이산 관찰 모델이며, 시간 오토마타 도구를 사용했다고 해서 물리적 시간 제약을 검증한 것은 아니다. Scenic과 VerifAI는 상황을 생성하고 시뮬레이션하여 자율 시스템을 시험한다~\cite{fremont2019scenic,dreossi2019verifai}. 본 연구는 환경을 탐색하지 않고 이미 기록되거나 작성된 사건 배열을 고정해 검사한다.

\subsection{본 연구의 위치}
\label{sec:related-position}

기존 연구는 한 장면의 설명과 행동의 정렬, 사람이 정한 규칙의 형식화와 검증, 개별 요구 문장의 형식화, 그리고 형식 명세에 대한 경로 감시를 다룬다. 본 연구는 이들 사이에 남은 질문, 즉 자동 생성된 연속 주석에서 어떤 문장이 의무를 만들고 어떤 증거가 그것을 해제하는지를 연결하고 증거가 없는 경우를 모순과 구분하는 문제를 다룬다. 형식 검증 기술을 새로 만들었다고 주장하지 않으며, 다른 방법이 이러한 연결을 할 수 없다고 단정하지도 않는다. 단순 FSM도 같은 규칙을 실행할 수 있으며, UPPAAL은 이 규칙을 질의로 표현하고 진단 경로를 남기는 구현 수단이다.

선행 기술과의 구체적인 접점은 조건의 대응 관계다. 한 사건에 행동을 정당화하는 조건 B와 의무를 해제하는 조건 A가 함께 언급되면, 구절을 역할별로 분리해 같은 조건 ID의 증거만 해당 의무에 적용한다. 모델 질의의 결과는 그 대응을 전제로 한 실행 결과다. 이 구조는 ``앞뒤 행동이 다르다''는 표시 대신 이전 요구가 아직 남았는지와 어떤 근거가 빠졌는지를 다시 읽게 한다. 표~\ref{tab:related-position}은 각 연구군의 입력ㆍ표현ㆍ출력과 본 논문의 범위를 정리한다.

\begin{table*}[t]
\centering
\color{revisionblue}
\bitablecaption{관련 방법의 입력ㆍ표현ㆍ출력과 본 논문의 범위.}{Inputs, representations, outputs, and scope relative to prior methods.}
\label{tab:related-position}
\scriptsize
\begin{tabularx}{\textwidth}{L{.21\textwidth}L{.24\textwidth}L{.24\textwidth}Y}
\toprule
방법 & 입력과 표현 & 출력 & 본 연구와의 관계 \\
\midrule
추론 기반 주행 \cite{sima2024drivelm,tian2024drivevlm,wang2025alpamayo} & 장면 관측, 언어 추론과 계획 & 설명ㆍ행동/계획 & 저장된 설명 사이의 의무 해제 연결은 본 연구의 검사 대상 \\
추론 충실성/CoT--행동 검증 \cite{turpin2023,lanham2023,yu2026drivethoughts,mayumu2026faithfulness,nguyen2026vladrivebench} & 생성 설명ㆍ궤적, 관측 및 입력/CoT 개입 & 충실성, 장면ㆍ행동 정렬, 개입 반응 & 한 장면 단위 평가; 본 논문은 저장된 연속 주석의 조건 연결 검사이며 인과 충실성은 미검증 \\
교통 규칙 형식화ㆍ검증 \cite{esterle2020,maierhofer2020,wan2024} & 사람이 정한 규칙, 시간 논리 & 규칙 준수, 규칙의 일관성ㆍ정확성 & 본 논문은 생성 주석에 암묵적으로 담긴 의무를 다룸 \\
자연어--시간 논리 변환 \cite{chen2023nl2tl,cosler2023nl2spec} & 개별 요구 문장 & 형식 명세, 구절--식 대응 & 문장 사이의 의무--해제 연결은 본 연구의 대상; 범용 번역기는 제시하지 않음 \\
실행 중 검증ㆍLTL/TLTL 모니터 \cite{leucker2009runtime,bauer2011runtime} & 명세와 유한 관측 경로 & 참/거짓/미결정, 최초 확정 & 경로 검사ㆍ세 값ㆍ최초 위반은 선행 기술; 본 논문의 증거 미상과 의미가 같지는 않음 \\
UPPAAL \cite{behrmann2004uppaal} & 오토마타와 상태 속성 & 속성 판정, 진단 경로 & 조건별 의무 상태와 최초 문제 기록을 실행하는 도구 \\
Scenic/VerifAI \cite{fremont2019scenic,dreossi2019verifai} & 상황 명세와 시험 시스템 & 생성 상황, 분석 결과 & 본 논문은 환경ㆍ제어 탐색 대신 고정 사건 배열을 재생 \\
본 논문 & CoC 구절, 의무ㆍ조건 ID, 작성된 증거 & 모순/근거 부족/양립, 최초 사건ㆍ의무ㆍ사유 & 원문과 상태 판정의 연결을 실증; 자동 NLPㆍ자연 정확도 미검증 \\
\bottomrule
\end{tabularx}
\end{table*}
\end{revision}
